\subsection{Binary determinant rays: a bounded-hole subfamily}

The complete zero classification for invariants of a single binary form
already rules out one possible source of unbounded realized holes.
The classification used here is known: Brouwer--Draisma--Popoviciu,
\emph{The degrees of a system of parameters of the ring of invariants
of a binary form}, \emph{Transformation Groups} 20 (2015),
953--967, Proposition~4.2, printed p.~957. The following is its
specialization to the canonical ordinary-plethysm grading, rather
than a new classification.

\begin{proposition}[P-F25: the known binary-form zero classification on stretching rays]
\label{prop:binarydeterminantzeros}
Let $d\ge2$, set $g=\gcd(d,2)$ and $q=2/g$, and take the ordinary
plethysm data
\[
 W=\operatorname{Sym}^d\mathbb C^2,\qquad
 \lambda=(q),\qquad \pi=(d/g,d/g).
\]
Its canonical stretching function is
\[
 f_d(s)=\dim\operatorname{Hom}_{GL_2}
   \bigl((\det)^{ds/g},\operatorname{Sym}^{qs}W\bigr).
 \tag{B-H1}
\]
For $d=2,3$, it is nonzero exactly when $s$ is even.
For $d\ge4$, its positive-integer zeros are exactly those in the table:
\[
\begin{array}{c|l}
 d&\{s\ge1:f_d(s)=0\}\\ \hline
 5&\{1,3,5,7\}\\
 6&\{1,3,5,7,9,11,13\}\\
 7&\{1,3,5\}\\
 10&\{1,3,5,7\}\\
 14&\{1,3,5\}\\
 9,11,13&\{1,3\}\\
 d\equiv2\pmod4,\ d\ge18&\{1,3\}\\
 \text{all other }d\ge4&\{1\}.
\end{array}
 \tag{B-H2}
\]
In particular, after excluding identically zero congruence classes,
every realized stretching hole has index at most thirteen. This
bound is attained by binary sextics, where $f_6(13)=0$ and
$f_6(s)>0$ for every $s\ge14$.
\end{proposition}

\begin{proof}
Write $R_d=(\operatorname{Sym}W)^{SL_2}$ and denote its ordinary
coefficient degree by $D$. The scalar $-I_2$ forces $dD$ to be
even whenever $(R_d)_D\ne0$, equivalently $q\mid D$.
At degree $D=qs$, a general linear transformation acts on this
invariant subspace by $(\det)^{dD/2}=(\det)^{ds/g}$.
Indeed it is a scalar times a special linear matrix, and the scalar
acts with weight $dD$. Evaluation at compatible powers of a volume
vector therefore identifies the canonical algebra in (B-H1) with
the whole ring $R_d$, regraded by $\deg (R_d)_{qs}=s$.
As in P-F19, this evaluation respects Cartan multiplication.

Consequently $f_d(s)$ is the invariant dimension denoted
$h^d_{qs}$ in the cited Proposition~4.2. Substituting $D=qs$ in
that proposition gives (B-H2). In particular, for odd $d$ all odd
ordinary degrees are absent before regrading, while their finite
even-degree exceptions are divided by two. For even $d$ the
ordinary and stretch gradings coincide. The two exceptional rings
of orders two and three have invariant degrees divisible by two
and four, respectively, giving the stated even-stretch support.
The maximum finite zero in the table is thirteen and occurs when
$d=6$. This proves all the asserted consequences of the known
classification.
\end{proof}

\paragraph{What this does not say about the source's strict SH.}
GCT VI, Definition~1.2.4, tests a nonzero constituent polynomial
$f_i(X)$ at \emph{every} integer $X\ge1$, not only those integers
in the residue class represented by $i$. Thus (B-H2), which
classifies the actual values $f_d(s)$, does not by itself prove
$s(f_d)\le13$, nor a bound for $p(f_d)$. Off-residue constituent
values are additional data. The bound above concerns realized
holes and excludes this family as a source of arbitrarily late
actual zero coefficients; a strict-SH or PH2 conclusion would need
an extra constituent-positivity argument.

\paragraph{Polynomial generator bounds and the remaining parameter gate.}
For this same binary family, precise polynomial bounds are already
known for different invariants. Derksen,
\emph{Polynomial bounds for rings of invariants},
\emph{Proceedings of the American Mathematical Society} 129 (2001),
955--963, Theorem~1.1, Proposition~1.2, and Example~2.1,
printed p.~956, gives
\[
 \sigma(\operatorname{Sym}^d\mathbb C^2)\le2d^3,
 \qquad
 \beta(\operatorname{Sym}^d\mathbb C^2)
       \le\frac32(d+1)d^6.
 \tag{B-H3}
\]
Here $\sigma$ bounds the degrees of a collection of homogeneous
invariants defining the null cone, and $\beta$ bounds algebra
generator degrees. The coordinate-ring convention in that source
agrees with $R_d$ because $\operatorname{Sym}^d\mathbb C^2$ is
self-dual as an $SL_2$-module. The source obtains the first inequality from
$SL_2\subset\mathbb A^4$, its one defining equation of degree two,
group dimension three, and matrix coefficients of the degree-$d$
representation of degree at most $d$. The representation has
finite kernel, as required by Proposition~1.2.

These statements do not supply a polynomial bound on the canonical
modular index. The same paper, Section~3, equation~(3.1), printed
p.~957, constructs a genuine homogeneous system of parameters by
raising null-cone equations to a common least-common-multiple
degree. Its Section~4 explicitly avoids this step when proving
the improved \emph{generator} bound: Proposition~4.2 concerns a
finite module over a possibly dependent collection of invariants,
and its proof places the parameters in a polynomial extension of
the original ring. It does not place a small-degree parameter
system in the original invariant ring.

For clarity, a verified explicit, though not polynomial, consequence
is available. Put
\[
 B_d=\left\lfloor\frac{2d^3}{q}\right\rfloor,
 \qquad L_d=\operatorname{lcm}(1,2,\ldots,B_d).
\]
In the stretch grading the null-cone equations can be chosen with
degrees at most $B_d$. Their powers of degree $L_d$ still define
the irrelevant maximal ideal set-theoretically. Homogeneous
normalization gives a parameter system with every degree $L_d$;
the prime-avoidance extension argument in P-F17 also bounds the
source's lexicographically minimal canonical index by $L_d$.
For $d\ge3$ the invariant ring has dimension $d-2$
(Brouwer--Draisma--Popoviciu, Section~1, p.~953, with the
$d\ge3$ scope made explicit in Proposition~2.1, p.~954).
The Cohen--Macaulay theorem and the global quasi-polynomial
input for the actual canonical ray, together with P-F9, yield
\[
 \chi(f_d)\le L_d,\qquad
 s(f_d)\le p(f_d)\le(d-3)L_d.
 \tag{B-H4}
\]
The case $d=2$ is the separately explicit polynomial ring with
stretch generator degree two and both indices zero. Formula
(B-H4) is a computable envelope, not a proof of polynomial PH2 or
PH3. Converting the already polynomial null-cone bounds into a
polynomial bound for actual parameter degrees remains a separate
family-specific gate.

The primary sources are
\url{https://pure.tue.nl/ws/files/17749340/Degrees2015.pdf}
and
\url{https://sites.lsa.umich.edu/hderksen/wp-content/uploads/sites/614/2018/05/A.I.a.7.pdf}.
